% Синтетический тестовый документ: базовые конструкции математической статьи.
% Реальных статей в корпусе нет и быть не должно (см. CLAUDE.md).
\documentclass[12pt]{article}
\usepackage[T2A]{fontenc}
\usepackage[utf8]{inputenc}
\usepackage[russian]{babel}
\usepackage{amsmath,amssymb,amsthm}
\usepackage{tikz}
\usepackage{localmacros}

\theoremstyle{plain}
\newtheorem{theorem}{Теорема}[section]
\newtheorem{lemma}[theorem]{Лемма}
\newtheorem*{conjecture}{Гипотеза}
\theoremstyle{definition}
\newtheorem{definition}[theorem]{Определение}
\theoremstyle{remark}
\newtheorem{remark}{Замечание}

\numberwithin{equation}{section}

\newcommand{\F}{\mathbb{F}}
\newcommand{\wt}[1]{\operatorname{wt}(#1)}
\DeclareMathOperator{\rank}{rank}

\title{Тестовый документ}
\author{Корпус tex2docx}

\begin{document}
\maketitle

\section{Введение}\label{sec:intro}

Пусть $\code{C} \subseteq \F_q^n$ "--- линейный код с порождающей
матрицей~$\matr{G}$, а $\vect{x}$ "--- вектор.

\begin{definition}\label{def:square}
Квадратом Шура--Адамара кода $\code{C}$ называется код
\begin{equation}\label{eq:square}
  \code{C}^{\star 2} = \langle \vect{c} \star \vect{c}' \mid
  \vect{c}, \vect{c}' \in \code{C} \rangle .
\end{equation}
\end{definition}

\begin{theorem}[Основная]\label{thm:main}
Для случайного кода $\code{C}$ размерности $k$ с большой вероятностью
\[
  \dim \code{C}^{\star 2} = \min\left\{n, \binom{k+1}{2}\right\}.
\]
\end{theorem}

\begin{proof}
Следует из~\eqref{eq:square} и леммы~\ref{lem:aux}.
\end{proof}

\begin{lemma}\label{lem:aux}
Если $\rank \matr{G} = k$, то $\wt{\vect{x}} \le n$.
\end{lemma}

\begin{conjecture}
Утверждение теоремы~\ref{thm:main} верно для любого $q$.
\end{conjecture}

\section{Выкладки}\label{sec:calc}

\begin{align}
  a &= b + c, \label{eq:first}\\
  d &= e \cdot f. \notag\\
  g &= h. \label{eq:third}
\end{align}

\begin{remark}
Формулы~\eqref{eq:first} и~\eqref{eq:third} пронумерованы, средняя "--- нет.
\end{remark}

\begin{figure}[ht]
  \centering
  \begin{tikzpicture}
    \draw[->] (0,0) -- (2,0) node[right] {$x$};
    \draw[->] (0,0) -- (0,2) node[above] {$y$};
    \draw (0,0) circle (1);
  \end{tikzpicture}
  \caption{Окружность}\label{fig:circle}
\end{figure}

На рис.~\ref{fig:circle} изображена окружность (см. раздел~\ref{sec:intro}).

\end{document}
